\rok{Фебруар 2026.}{В}

Други практични тест из Алгоритама и структура података

Први практични тест одржан 11.02.2026.

\zadatak{1}{Завршавање имплементације BFS алгоритма}
Основна метода BFS алгоритма није потпуна. Допунити имплементацију ове методе која треба да омогући комплетно извршавање овог алгоритма.

\textbf{Додатне напомене за решавање задатка:}
\begin{itemize}
	\item Написати алгоритам.
	\item У Main методи креирати произволан граф. Обезбедити начин за тестирање и проверу нове функционалности, односно позив BFS методе.
	\item У template-у је закоментарисана метода \texttt{public void BFS(int startVertex)}.
\end{itemize}

\zadatak{2}{Проналажење карактера са максималним бројем појављивања}
Написати методу која идентификује карактер који се највише пута појављује у string-у.

\textbf{Додатни захтеви за решавање задатка:}
\begin{itemize}
	\item Користити Hash табелу из шаблона за чување броја појављивања сваког карактера.
	\item Имплементацију писати у методи:
	\begin{Verbatim}[fontsize=\small]
public static void maxOccurrence(string str)
	\end{Verbatim}
	\item Разлика између малих и великих слова постоји.
	\item Уколико постоји више карактера који се појављују максималан број пута, исписати их све.
\end{itemize}

\textbf{Пример:}
\begin{itemize}
	\item \textbf{Улаз 1:} \texttt{str = "Algoritmi i strukture"}
	\item \textbf{Излаз 1:} \texttt{"i"} (појављује се 3 пута)
	\item \textbf{Улаз 2:} \texttt{str = "Vetar"}
	\item \textbf{Излаз 2:} \texttt{"Vetar"} (сви карактери се појављују само једном)
\end{itemize}

\zadatak{3}{Путања са наизменичном парношћу (Zig-Zag Парност)}
Написати методу која проверава да ли постоји путања од корена до листа на којој се парни и нечетни бројеви наизменично смењују (нпр. парни -> нечетни -> парни...).

\textbf{Додатни захтеви за решавање задатка:}
\begin{itemize}
	\item Јавна метода:
	\begin{Verbatim}[fontsize=\small]
public bool hasAlternatingPath()
	\end{Verbatim}
	\item Приватна метода:
	\begin{Verbatim}[fontsize=\small]
private bool hasAlternatingPath(Node current, bool expectedEven)
	\end{Verbatim}
	\item Подразумева се да стабло има барем два чвора.
	\item Алгоритам ОБАВЕЗНО писати у оквиру класе \texttt{BinarySearchTree}.
\end{itemize}

\textbf{Примери:}

\begin{center}
\begin{minipage}{\textwidth}
\centering
\begin{forest}
for tree={
	circle,
	draw,
	thick,
	fill=gray!20,
	minimum size=8mm,
	inner sep=1pt,
	s sep=8mm,
	l sep=12mm,
	edge={-, thick},
	font=\small
}
[8
	[4
		[2
			[, phantom]
			[3]
		]
		[5]
	]
	[9
		[, phantom]
		[11]
	]
]
\end{forest}

\vspace{0.3cm}
\textbf{Слика 1.} Пример BST-a.
\end{minipage}
\end{center}

\begin{itemize}
	\item \textbf{Улаз 1:} \texttt{/}
	\item \textbf{Излаз 1:} \texttt{false}
\end{itemize}

\begin{center}
\begin{minipage}{\textwidth}
\centering
\begin{forest}
for tree={
	circle,
	draw,
	thick,
	fill=gray!20,
	minimum size=8mm,
	inner sep=1pt,
	s sep=8mm,
	l sep=12mm,
	edge={-, thick},
	font=\small
}
[10
	[4
		[2
			[3]
			[, phantom]
		]
		[6
			[5]
			[9]
		]
	]
	[20
		[12
			[, phantom]
			[14]
		]
		[29
			[27]
			[, phantom]
		]
	]
]
\end{forest}

\vspace{0.3cm}
\textbf{Слика 2.} Пример BST-a.
\end{minipage}
\end{center}

\begin{itemize}
	\item \textbf{Улаз 2:} \texttt{/}
	\item \textbf{Излаз 2:} \texttt{false}
\end{itemize}

\sablon{Шаблон другог теста фебруар 2026. група В}{2025/Februar/GrupaC/AISP2025_Test2_GrupaC.linq}
